\chapter{अवकल कलन}\label{ch:b2:diffcalc}

प्रथम वर्ष के खंड का दो-चर कलन यहाँ मानदंडित समष्टियों के बीच प्रतिचित्रणों
के अवकल कलन में परिपक्व होता है: सर्वोत्तम रैखिक सन्निकटन के रूप में
\emph{\omterm{def:b2:diffcalc:differential}{अवकल}}, पूरी व्यापकता में शृंखला नियम, \emph{सिद्ध} की गई श्वार्ज़ की
सममिति प्रमेय, टेलर सूत्र, और चरम मानों का पूरा द्वितीय-कोटि विश्लेषण। और
सिद्धांत का मुकुट, प्रतिलोम फलन प्रमेय, अपनी उपपत्ति-रणनीति --- बानाख अचल
बिंदु --- को स्पष्ट करते हुए कहा गया है।

आगे सर्वत्र $U$ $\R^n$ (अथवा किसी मानदंडित समष्टि; परिमित-विमीय स्थिति सारे
विचार उठा लेती है) का कोई \omterm{def:b2:metric:topology}{विवृत} उपसमुच्चय है, $f \colon U \to
\R^m$।

\section{अवकल}

\begin{definition}\label{def:b2:diffcalc:differential}
$f$ $a$ पर \emph{अवकलनीय}\index{अवकलनीय} कहलाता है जब कोई (\omterm{def:b2:metric:continuity}{संतत}) रैखिक
प्रतिचित्रण $\dd f_a \colon \R^n \to \R^m$ ऐसा हो कि
\[
f(a + h) = f(a) + \dd f_a(h) + o(\norm h)
\qquad (h \to 0).
\]
प्रतिचित्रण $\dd f_a$, अर्थात् $a$ पर $f$ का \emph{अवकल}\index{अवकल}, अद्वितीय
होता है; और विहित आधारों में उसका आव्यूह \emph{याकोबी आव्यूह}\index{याकोबी
आव्यूह} $J_f(a) =
\bigl(\frac{\partial f_i}{\partial x_j}(a)\bigr)$ है। अवकलनीयता से संततता तथा सभी दिशिक अवकलजों $\dd f_a(v) = \lim_{t\to0}\frac{f(a + tv)
- f(a)}{t}$ का अस्तित्व
निकलता है; पर विलोम विफल है (\cref{exo:b2:diffcalc:2})। $m = 1$ के लिए $\dd f_a(h) = \langle
\nabla f(a), h\rangle$: अर्थात् प्रथम वर्ष
का प्रवणक, अब अवकल को निरूपित करने वाले सदिश के रूप में समझा हुआ।
\end{definition}

\begin{theorem}[$C^1$ कसौटी]\label{thm:b2:diffcalc:c1}
यदि $f$ के सारे आंशिक अवकलज $U$ पर विद्यमान हों और $a$ पर \omterm{def:b2:metric:continuity}{संतत} हों, तो
$f$ $a$ पर \omterm{def:b2:diffcalc:differential}{अवकलनीय} है। अतः ``$U$ पर $C^1$'' --- अर्थात् \omterm{def:b2:metric:continuity}{संतत} आंशिक अवकलज ---
से सर्वत्र अवकलनीयता निकलती है, \omterm{def:b2:metric:continuity}{संतत} \omterm{def:b2:diffcalc:differential}{अवकल} के साथ।
\end{theorem}

\begin{proof}
घटक-दर-घटक ($m = 1$ पर्याप्त है)। दो चरों के लिए प्रथम वर्ष की उपपत्ति ---
एक-एक निर्देशांक बदलिए, हर पैर पर एक-चर मध्यमान प्रमेय लगाइए, और $a$ पर
आंशिक अवकलजों की संततता का उपयोग कीजिए --- $n$ पैरों तक अक्षरशः सामान्यीकृत
हो जाती है:
\[
f(a + h) - f(a) = \sum_{j=1}^{n} \bigl(f(a + h^{(j)}) - f(a +
h^{(j-1)})\bigr)
= \sum_j h_j\,\frac{\partial f}{\partial x_j}(\xi_j),
\]
जहाँ $h^{(j)}$ $h$ के पहले $j$ निर्देशांक जमा देता है और $\xi_j$ $j$-वें पैर पर
पड़ता है; और संततता हर $\frac{\partial f}{\partial x_j}(\xi_j)$ को $\frac{\partial
f}{\partial x_j}(a) + o(1)$ बना देती है, तथा त्रुटि $o(\norm h)$ रह जाती
है।
\end{proof}

\begin{theorem}[शृंखला नियम]\label{thm:b2:diffcalc:chainrule}
यदि $f$ $a$ पर \omterm{def:b2:diffcalc:differential}{अवकलनीय} हो और $g$ $f(a)$ पर, तो $g \circ f$ $a$ पर \omterm{def:b2:diffcalc:differential}{अवकलनीय} है और
\[
\dd(g \circ f)_a = \dd g_{f(a)} \circ \dd f_a ,
\qquad
J_{g\circ f}(a) = J_g\bigl(f(a)\bigr)\,J_f(a) :
\]
अर्थात् याकोबी गुणा हो जाते हैं।\index{शृंखला नियम}
\end{theorem}

\begin{proof}
$f(a + h) = f(a) + \dd f_a(h) + \norm h\,\varepsilon_1(h)$ और $b = f(a)$, $k = k(h) = \dd f_a(h) + \norm h \varepsilon_1(h)$ के साथ $g(b + k) = g(b) + \dd g_b(k) + \norm k\,\varepsilon_2(k)$ लिखिए। प्रतिस्थापित करने पर,
\[
g(f(a+h)) = g(b) + \dd g_b\bigl(\dd f_a(h)\bigr)
+ \norm h\,\dd g_b(\varepsilon_1(h)) + \norm{k}\,\varepsilon_2(k),
\]
और दोनों त्रुटि-पद $o(\norm h)$ हैं: पहला इसलिए कि $\dd g_b$ \omterm{def:b2:metric:continuity}{संतत} है और $\varepsilon_1 \to 0$; और दूसरा
इसलिए कि $\norm k
\leq C\norm h$ (परिबद्ध रैखिक प्रतिचित्रण जमा छोटा पद) और $h \to 0$ होने पर $\varepsilon_2(k) \to 0$।
\end{proof}

\begin{example}[त्रिज्यीय फलन, एक बार में सदा के लिए]\label{ex:b2:diffcalc:radial}
$\R^n\setminus\{0\}$ पर $r(x) = \norm x_2$ लीजिए और $f = g \circ
r$ लीजिए, जहाँ $g$ एक चर का $C^1$ फलन है। पहले, $r$
$0$ से दूर \omterm{def:b2:diffcalc:differential}{अवकलनीय} है: $r^2 = \sum x_i^2$ से
\[
\frac{\partial r}{\partial x_i} = \frac{x_i}{r},
\qquad\text{अर्थात्}\qquad
\nabla r(x) = \frac{x}{\norm x} ,
\]
अर्थात् इकाई त्रिज्यीय सदिश ($r^2$ का अवकलन करके भाग दीजिए --- अथवा $\sqrt{\cdot}$
पर शृंखला नियम लगाइए)। फिर शृंखला नियम हर त्रिज्यीय फलन के लिए देता है
\[
\nabla f(x) = g'\bigl(\norm x\bigr)\,\frac{x}{\norm x} .
\]
किया हुआ उदाहरण: $g(r) = \frac1r$ $\nabla\frac{1}{\norm
x} = -\frac{x}{\norm x^3}$ देता है, अर्थात् गुरुत्वाकर्षण और
स्थिरवैद्युतिकी का प्रतिलोम-वर्ग क्षेत्र --- दिशा त्रिज्यीय, परिमाण $\frac{1}{\norm x^2}$।
समापन दृष्टि: त्रिज्यीय फलनों के प्रवणक इसलिए त्रिज्यीय होते हैं कि उनके
स्तर समुच्चय गोले हैं और प्रवणक स्तर समुच्चयों के लांबिक होता है; इसके
विपरीत $x = 0$ पर $r$ \omterm{def:b2:diffcalc:differential}{अवकलनीय} \emph{नहीं} है (कोई भी उम्मीदवार रैखिक
प्रतिचित्रण सभी दिशाओं से $\norm h$ से मेल नहीं खाता) --- मूल बिंदु को शालीनता
से पार करने के लिए त्रिज्यीय परिच्छेदिकाओं को $g'(0) = 0$ चाहिए।
\end{example}

\begin{theorem}[मध्यमान असमिका]\label{thm:b2:diffcalc:mvi}
मान लीजिए $f$ $U$ पर \omterm{def:b2:diffcalc:differential}{अवकलनीय} है और खंड $\intcc{a}{b}
= \{a + t(b-a)\}$ $U$ में पड़ता है। तब
\[
\norm{f(b) - f(a)} \leq \norm{b - a}\,
\sup_{x \in \intcc{a}{b}} \vertiii{\dd f_x} .
\]
विशेष रूप से, किसी \emph{\omterm{def:b2:metric:connected}{संबद्ध}} \omterm{def:b2:metric:topology}{विवृत} समुच्चय पर शून्य \omterm{def:b2:diffcalc:differential}{अवकल} वाला
\omterm{def:b2:diffcalc:differential}{अवकलनीय} प्रतिचित्रण अचर होता है।
\end{theorem}

\begin{proof}
फलन $\varphi(t) = f(a + t(b-a))$ $\intcc{0}{1}$ पर \omterm{def:b2:diffcalc:differential}{अवकलनीय} है और $\varphi'(t) = \dd f_{a + t(b-a)}(b - a)$ (शृंखला नियम), जिसका \omterm{def:b2:nvs:norm}{मानदंड} $\leq M\norm{b-a}$ है जहाँ
$M$ वही प्रदर्शित उच्चतम है। $\R$-मान वाले $f$ के लिए एक-चर मध्यमान असमिका
निष्कर्ष दे देती है; और सदिश मानों के लिए उसे $t \mapsto \langle u,
\varphi(t)\rangle$ पर लगाइए, जहाँ $u$ $f(b) - f(a)$
के अनुदिश इकाई सदिश है। अचरता: स्थानीय अचरता (गोलों के भीतर खंड) और
संबद्धता (जिस समुच्चय पर $f$ किसी दिए हुए मान के बराबर है वह \omterm{def:b2:metric:topology}{विवृत} भी है
और संवृत भी: \cref{ch:b2:metric})।
\end{proof}

\begin{example}[मध्यमान असमिका से एक लिप्शिट्स अचर]\label{ex:b2:diffcalc:lipschitz}
क्या $f(x, y) = \sin x\,\sin y$ $\R^2$ पर \omterm{def:b2:metric:continuity}{लिप्शिट्स} है, और किस अचर के साथ? उसका प्रवणक $\nabla f = (\cos x\sin y,\
\sin x\cos y)$ है,
जिसका वर्ग-मानदंड
\[
\cos^2x\sin^2y + \sin^2x\cos^2y
\leq \sin^2 y + \cos^2y\cdot 1 = 1
\]
है ($\cos^2x$ और $\sin^2x$ को अलग-अलग $1$ से परिबद्ध कीजिए), अतः सर्वत्र $\vertiii{\dd f_{(x,y)}} = \norm{\nabla f} \leq 1$; और किन्हीं
दो बिंदुओं के बीच के खंड पर \cref{thm:b2:diffcalc:mvi} देता है
\[
\abs{f(b) - f(a)} \leq \norm{b - a}_2 :
\]
अर्थात् $f$ $1$-लिप्शिट्स है, और अचर तीखा है (मूल बिंदु के पास $f(x, \tfrac\pi2) = \sin x$ की
ढाल $1$ है)। समापन दृष्टि: मध्यमान असमिका \omterm{def:b2:diffcalc:differential}{अवकल} के किसी \omterm{def:b2:funcseq:def}{बिंदुवार} परिबंध
को संततता के वैश्विक मापांक में बदल देती है --- और हर विमा में \omterm{def:b2:metric:continuity}{लिप्शिट्स}
आकलनों का यही मानक रास्ता है, तथा \cref{exo:b2:diffcalc:12} के भीतर का इंजन भी।
\end{example}

\section{द्वितीय अवकलज}

\begin{theorem}[श्वार्ज़]\label{thm:b2:diffcalc:schwarz}
यदि $f$ $U$ पर $C^2$ हो (अर्थात् सारे द्वितीय आंशिक अवकलज विद्यमान और \omterm{def:b2:metric:continuity}{संतत}
हों), तो सभी $i, j$ के लिए:\index{श्वार्ज़ की प्रमेय}
\[
\frac{\partial^2 f}{\partial x_i\,\partial x_j}
= \frac{\partial^2 f}{\partial x_j\,\partial x_i} .
\]
\end{theorem}

\begin{proof}
दो चर पर्याप्त हैं ($x = x_i$, $y = x_j$, शेष जमे हुए)। द्वितीय अंतर
\[
\Delta(h) = f(a + h, b + h) - f(a + h, b) - f(a, b + h) + f(a,b) .
\]
लीजिए। $h$ स्थिर कीजिए और $\varphi(x) = f(x, b+h) - f(x, b)$ रखिए: तब $\Delta(h)
= \varphi(a + h) - \varphi(a)$, और मध्यमान प्रमेय के दो प्रयोग
\[
\Delta(h) = h\,\varphi'(\xi)
= h\Bigl(\frac{\partial f}{\partial x}(\xi, b+h) -
\frac{\partial f}{\partial x}(\xi, b)\Bigr)
= h^2\,\frac{\partial^2 f}{\partial y\,\partial x}(\xi, \eta),
\]
देते हैं, जहाँ $\xi \in \intoo{a}{a+h}$, $\eta \in \intoo{b}{b+h}$। संततता से $h \to 0$ होने पर $\frac{\Delta(h)}{h^2} \to \frac{\partial^2 f}{\partial
y\partial x}(a, b)$। और चरों की भूमिकाएँ
बदलकर वही परिकलन (पहले दूसरा चर जमाइए) $\frac{\Delta(h)}{h^2} \to \frac{\partial^2 f}{\partial x
\partial y}(a,b)$ देता है: अर्थात् एक ही राशि
की दोनों सीमाएँ मेल खाती हैं।
\end{proof}

\begin{example}[$C^2$ क्यों चाहिए: पियानो का प्रतिउदाहरण]\label{ex:b2:diffcalc:peano}
$f(x, y) = \dfrac{xy(x^2 - y^2)}{x^2 + y^2}$, $f(0,0) = 0$ लीजिए। मूल बिंदु से दूर $f$ $C^\infty$ है; और मूल बिंदु पर सारे प्रथम
तथा द्वितीय आंशिक अवकलज विद्यमान हैं, पर मिश्रित अवकलज असहमत हैं। अक्षों
के अनुदिश परिकलन कीजिए: $f(x, 0) = f(0, y) = 0$, और $y \neq 0$ के लिए
\[
\frac{\partial f}{\partial x}(0, y)
= \lim_{x\to0}\frac{f(x,y)}{x}
= \frac{y(0 - y^2)}{y^2} = -y ,
\qquad\text{सममित रूप से}\qquad
\frac{\partial f}{\partial y}(x, 0) = x .
\]
अतः
\[
\frac{\partial^2 f}{\partial y\,\partial x}(0,0)
= \frac{\dd}{\dd y}\Bigl[\frac{\partial f}{\partial
x}(0,y)\Bigr]_{y=0} = -1,
\qquad
\frac{\partial^2 f}{\partial x\,\partial y}(0,0) = +1 :
\]
अर्थात् दोनों मिश्रित आंशिक अवकलज विद्यमान हैं और भिन्न हैं। \cref{thm:b2:diffcalc:schwarz} के साथ
कोई विरोधाभास नहीं: $f$ के द्वितीय आंशिक अवकलज $0$ पर \omterm{def:b2:metric:continuity}{संतत} नहीं हैं ($y = tx$
के अनुदिश जाँचिए)। समापन दृष्टि: श्वार्ज़ की प्रमेय \emph{संततता} के बारे
में सच्ची प्रमेय है, कोई औपचारिक सर्वसमिका नहीं --- और नीचे टेलर--यंग में
परिकल्पना ``$C^2$'' असली काम कर रही है।
\end{example}

\begin{theorem}[कोटि 2 पर टेलर--यंग]\label{thm:b2:diffcalc:taylor}
मान लीजिए $f \colon U \to \R$ $C^2$ है और $a \in U$। तब $h \to 0$ होने पर
\[
f(a + h) = f(a) + \langle\nabla f(a), h\rangle
+ \frac12\, \langle H_a h,\, h\rangle + o\bigl(\norm h^2\bigr),
\]
जहाँ $H_a = \bigl(\frac{\partial^2 f}{\partial x_i\partial
x_j}(a)\bigr)$ \emph{हेस्सियन आव्यूह}\index{हेस्सियन आव्यूह} है (जो श्वार्ज़ से
\omterm{def:b2:quadratic:adjoint}{सममित} है)।
\end{theorem}

\begin{proof}
$\intcc{0}{1}$ पर $\varphi(t) = f(a + th)$ पर एक-चर टेलर--यंग प्रमेय (प्रथम वर्ष का खंड) लगाइए: शृंखला
नियम से $\varphi'(t) = \langle \nabla f(a + th), h\rangle$ और $\varphi''(t)
= \langle H_{a+th}h, h\rangle$, और दोनों $t$ में \omterm{def:b2:metric:continuity}{संतत} हैं। तब $\theta \in \intoo{0}{1}$ के साथ
$\varphi(1) = \varphi(0) + \varphi'(0) + \frac12\varphi''(\theta)$
(टेलर--लाग्रांज), और द्वितीय आंशिक अवकलजों की संततता $\varphi''(\theta) =
\varphi''(0) + o(1)\cdot\norm h^2$-एकसमान रूप में
बदल देती है: यही प्रदर्शित प्रसार है।
\end{proof}

\begin{example}[एक प्रसार, दो तरह से]\label{ex:b2:diffcalc:taylortwoways}
$f(x, y) = \eu^x\cos y$ का मूल बिंदु पर कोटि $2$ तक प्रसार कीजिए। \emph{एक-चर प्रसारों के
संयोजन से:}
\[
\eu^x\cos y = \Bigl(1 + x + \frac{x^2}{2} +
o(x^2)\Bigr)\Bigl(1 - \frac{y^2}{2} + o(y^2)\Bigr)
= 1 + x + \frac{x^2 - y^2}{2} + o\bigl(\norm{(x,y)}^2\bigr) .
\]
\emph{आंशिक अवकलजों से:} $f_x = \eu^x\cos y$, $f_y =
-\eu^x\sin y$, अतः $\nabla f(0) = (1, 0)$; और $f_{xx} = f$, $f_{yy} = -f$, $f_{xy} = -\eu^x\sin y$ से $H_0 =
\operatorname{diag}(1, -1)$:
अर्थात् \cref{thm:b2:diffcalc:taylor} $1 + x + \frac12(x^2 - y^2)$ को पुनः उत्पन्न कर देता है। दोनों परिकलन सहमत हैं, और
संयोजन वाला रास्ता तेज़ था --- कोई द्वितीय आंशिक अवकलज लगा ही नहीं। समापन
दृष्टि: मूल बिंदु क्रांतिक बिंदु \emph{नहीं} है ($\nabla f \neq 0$), अतः अनिश्चित
हेस्सियन के बावजूद यहाँ कोई काठी घोषित करने को नहीं है: रैखिक पद ही शासन
करता है, और द्वितीय-कोटि जाँच केवल क्रांतिक बिंदुओं पर ही बोलती है।
\end{example}

\begin{theorem}[द्वितीय-कोटि चरम-मान जाँच, सिद्ध]\label{thm:b2:diffcalc:extrema}
मान लीजिए $f$ किसी क्रांतिक बिंदु $a$ ($\nabla f(a) = 0$) के पास $C^2$ है, और उसका
हेस्सियन $H = H_a$ है।
\begin{enumerate}
  \item यदि $H$ धनात्मक निश्चित हो, तो $a$ यथार्थ स्थानीय न्यूनतम है
        (ऋणात्मक निश्चित: महत्तम)।
  \item यदि $H$ के \omterm{def:b2:reduction:eigen}{अभिलक्षणिक मान} दोनों चिह्नों के हों, तो $a$ काठी है:
        कोई चरम मान नहीं।
  \item यदि $H$ अपभ्रष्ट (और अर्धनिश्चित) हो, तो कोई निष्कर्ष नहीं।
\end{enumerate}
प्रथम वर्ष की ``$rt - s^2$'' जाँच $n = 2$ वाली स्थिति है: $\det H = rt -
s^2$, और $r$ से $\operatorname{tr}$-चिह्न
पढ़ लिया जाता है।
\end{theorem}

\begin{proof}
वर्णक्रमीय प्रमेय (\cref{thm:b2:quadratic:spectral}) से $H$ का \omterm{def:b2:quadratic:def}{द्विघात रूप} उसके चरम अभिलक्षणिक मानों
के बीच निचुड़ जाता है: $\lambda_{\min}\norm h^2 \leq \langle Hh, h\rangle \leq
\lambda_{\max}\norm h^2$।

(1) यदि $\lambda_{\min} > 0$: तो टेलर--यंग छोटे $h \neq 0$ के लिए
\[
f(a + h) - f(a) \geq \frac{\lambda_{\min}}{2}\norm h^2 -
o(\norm h^2) > 0
\]
देता है: अर्थात् यथार्थ स्थानीय न्यूनतम।

(2) $\lambda_+ > 0$ वाले किसी \omterm{def:b2:reduction:eigen}{अभिलक्षणिक सदिश} $v_+$ के अनुदिश: छोटे $t$ के लिए $f(a + tv_+) -
f(a) = \frac{\lambda_+}{2}t^2 + o(t^2) > 0$; और
$\lambda_- < 0$ वाले $v_-$ के अनुदिश अंतर ऋणात्मक है: अतः हर प्रतिवेश में दोनों चिह्न
आते हैं।

(3) $f(x,y) = x^2 + y^4$, $x^2 - y^4$, $x^2 + y^3$ का $0$ पर वही अपभ्रष्ट अर्धनिश्चित हेस्सियन है, पर तीन
भिन्न व्यवहार।
\end{proof}

\begin{example}[एक पूरा वर्गीकरण, वैश्विक सहित]\label{ex:b2:diffcalc:quarticrun}
$\R^2$ पर $f(x, y) = x^4 + y^4 - 4xy$ के सारे चरम मानों का वर्गीकरण कीजिए। \emph{क्रांतिक बिंदु:}
$\nabla f = (4x^3 - 4y,\ 4y^3 - 4x) =
0$ से $y = x^3$ और $x = y^3 = x^9$ मिलते हैं, अतः $x(x^8 - 1) = 0$: और वास्तविक हल $(0,0)$, $(1,1)$, $(-1,-1)$ हैं।
\emph{हेस्सियन:} $H = \begin{pmatrix} 12x^2 & -4\\ -4 &
12y^2\end{pmatrix}$। $(\pm1, \pm1)$ पर: $\begin{pmatrix} 12 &
-4\\ -4 & 12\end{pmatrix}$, \omterm{def:b2:reduction:eigen}{अभिलक्षणिक मान} $8$ और $16$:
अर्थात् धनात्मक निश्चित, यानी $f = -2$ के साथ यथार्थ स्थानीय न्यूनतम। $(0,0)$
पर: $\begin{pmatrix} 0 & -4\\ -4 & 0\end{pmatrix}$, \omterm{def:b2:reduction:eigen}{अभिलक्षणिक मान} $\pm4$: अर्थात् काठी। \emph{वैश्विकता:} $2\abs{xy} \leq x^2 +
y^2$ से
\[
f(x, y) \geq x^4 + y^4 - 2(x^2 + y^2)
= (x^2 - 1)^2 + (y^2 - 1)^2 + x^2 + y^2 - 2
\xrightarrow[\norm{(x,y)}\to\infty]{} +\infty :
\]
अर्थात् $f$ प्रबल है, अतः वह वैश्विक न्यूनतम प्राप्त करता है
(उपस्तर-समुच्चयों की संहतता), और अनिवार्यतः किसी क्रांतिक बिंदु पर: इसलिए
मान $-2$, जो $(1,1)$ तथा $(-1,-1)$ दोनों पर है, वैश्विक न्यूनतम है; और कोई महत्तम
नहीं है ($f$ ऊपर से अपरिबद्ध)। समापन दृष्टि: स्थानीय जाँच उम्मीदवारों का
वर्गीकरण करती है, पर ``स्थानीय'' को ``वैश्विक'' केवल कोई वृद्धि-तर्क बनाता
है --- और यही इस पुस्तक की हर अनुकूलन-उपपत्ति का दो-चरणी ढर्रा है।
\end{example}

\begin{method}[$f \colon \R^n \to \R$ के चरम मानों का वर्गीकरण]\label{met:b2:diffcalc:extrema}
\begin{enumerate}
  \item $\nabla f = 0$ हल कीजिए (सारे क्रांतिक बिंदु; और सीमा वाले प्रांत पर सीमा
        को अलग से निपटाइए, जैसा \cref{exo:b2:diffcalc:7} में)।
  \item हर क्रांतिक बिंदु पर हेस्सियन और उसके अभिलक्षणिक मानों के चिह्न
        परिकलित कीजिए --- विमा $2$ में तो बस $\det H$ और $\operatorname{tr} H$: $\det < 0$ काठी; $\det >
        0$
        चरम मान, जिसका प्रकार अनुरेख के चिह्न से मिलता है; और $\det = 0$: जाँच
        चुप है, वक्रों के अनुदिश $f$ का अध्ययन कीजिए।
  \item वैश्विक कथनों के लिए कोई संहतता या प्रबलता वाला तर्क जोड़िए (अनंत
        पर $f \to +\infty$, अथवा कोई \omterm{def:b2:metric:compact}{संहत} प्रतिबंध-समुच्चय), और फिर क्रांतिक मानों की
        तुलना कीजिए।
\end{enumerate}
\end{method}

\section{प्रतिलोम फलन प्रमेय}

\begin{theorem}[प्रतिलोम फलन प्रमेय]\label{thm:b2:diffcalc:inverse}
मान लीजिए $f \colon U \to \R^n$ $C^1$ है और $a \in U$, जहाँ $\dd f_a$ \emph{प्रतिलोमनीय} है। तब ऐसे \omterm{def:b2:metric:topology}{विवृत}
प्रतिवेश $V \ni a$, $W
\ni f(a)$ हैं कि $f \colon V \to W$ एकैकी आच्छादक हो और उसका प्रतिलोम $C^1$ हो, तथा
\index{प्रतिलोम फलन प्रमेय}
\[
\dd (f^{-1})_{f(x)} = (\dd f_x)^{-1} \qquad (x \in V).
\]
\end{theorem}
\admitted

\begin{remark}[यह सत्य क्यों है: अचल-बिंदु रणनीति]
$a$ के पास $f(x) = y$ हल करना अचल-बिंदु समीकरण $x = x + \dd f_a^{-1}\bigl(y - f(x)\bigr) =: \Phi_y(x)$ के रूप में लिखा जा सकता है;
और प्रतिचित्रण $\Phi_y$ का \omterm{def:b2:diffcalc:differential}{अवकल} $\mathrm{id} - \dd f_a^{-1}\dd f_x$ है, जो $\dd f$ की संततता से $a$ के पास छोटा
है, अतः $\Phi_y$ किसी छोटे संवृत गोले पर संकुचन है और बानाख अचल बिंदु प्रमेय
(\cref{thm:b2:metric:banach}) अद्वितीय स्थानीय हल $x = f^{-1}(y)$ दे देती है। और फिर प्रतिलोम की संततता तथा
अवकलनीयता संकुचन के आकलनों से निकल आती हैं। पूरा लेखा-जोखा तृतीय वर्ष में
किया गया है; रणनीति --- और कथन --- अब से खुलकर प्रयुक्त होते हैं। और साथी
\emph{अंतर्निहित फलन प्रमेय} ($\frac{\partial F}{\partial y}$ के प्रतिलोमनीय होने पर $y(x)$ के लिए $F(x, y) = 0$
हल करना) $(x, y) \mapsto (x, F(x,y))$ पर प्रमेय लगाने से निकल आती है।
\end{remark}

\begin{example}[ध्रुवीय निर्देशांक]\label{ex:b2:diffcalc:polar}
$\Phi(r, \theta) = (r\cos\theta, r\sin\theta)$ का याकोबी
\[
J_\Phi = \begin{pmatrix} \cos\theta & -r\sin\theta\\ \sin\theta &
r\cos\theta \end{pmatrix},
\qquad \det J_\Phi = r :
\]
है, जो $r \neq 0$ के लिए प्रतिलोमनीय है, अतः $\Phi$ मूल बिंदु से दूर स्थानीय $C^1$
अवकल समाकृतिकता है --- अर्थात् ``ध्रुवीय निर्देशांकों में बदलने'' का
लाइसेंस, जो \cref{ch:b2:multint} के बहु समाकलों के लिए फिर से नवीनीकृत हुआ।
\end{example}

\begin{example}[सर्वत्र स्थानीय, कहीं भी वैश्विक नहीं]\label{ex:b2:diffcalc:localnotglobal}
$\R^2$ पर $f(x, y) = \bigl(\eu^x\cos y,\ \eu^x\sin y\bigr)$ लीजिए। उसका याकोबी,
\[
J_f = \begin{pmatrix}
\eu^x\cos y & -\eu^x\sin y\\
\eu^x\sin y & \eu^x\cos y
\end{pmatrix},
\qquad
\det J_f = \eu^{2x} > 0 ,
\]
कभी लुप्त नहीं होता: अतः \cref{thm:b2:diffcalc:inverse} के अनुसार $f$ समतल के \emph{प्रत्येक} बिंदु
पर स्थानीय $C^1$ अवकल समाकृतिकता है। फिर भी $f$ एकैकी होने से बहुत दूर है:
$f(x, y + 2\pi) = f(x, y)$, अतः हर मान अपरिमित बार लिया जाता है; और वह आच्छादक भी नहीं है,
क्योंकि $\norm{f(x, y)} = \eu^x > 0$ मूल बिंदु छोड़ देता है। समापन दृष्टि: प्रतिलोम फलन प्रमेय
अखंडनीय रूप से \emph{स्थानीय} है --- हर $\dd f_a$ की प्रतिलोमनीयता स्थानीय
प्रतिलोमों की एक चकत्ती देती है, और उनका किसी एक प्रतिलोम में जुड़ जाना
आवश्यक नहीं। (जो पाठक सम्मिश्र संख्याएँ जानते हैं वे $z \mapsto
\eu^z$ पहचान लेंगे; और
यह चकत्ती लघुगणक की शाखाओं का कुल है।) इसकी तुलना \cref{exo:b2:diffcalc:12} से कीजिए, जहाँ कोई
परिमाणात्मक वैश्विक परिकल्पना सचमुच एक ही वैश्विक प्रतिलोम को बाध्य कर
देती है।
\end{example}

\begin{remark}[सामान्य भूलें]
\emph{(क) दिशिक अवकलज सस्ते हैं, \omterm{def:b2:diffcalc:differential}{अवकल} नहीं:} सारे दिशिक अवकलज विद्यमान हो
सकते हैं --- और दिशा पर रैखिक रूप से निर्भर न भी हों --- फिर भी अवकलनीयता
न हो (\cref{exo:b2:diffcalc:2}); केवल $C^1$ कसौटी (\cref{thm:b2:diffcalc:c1}) आंशिक अवकलजों को \omterm{def:b2:diffcalc:differential}{अवकल} तक बढ़ाती है।
\emph{(ख) क्रांतिक का अर्थ चरम नहीं:} काठियाँ (\cref{ex:b2:diffcalc:quarticrun}) और चुप अपभ्रष्ट स्थिति
(\cref{thm:b2:diffcalc:extrema}~(3)) दोनों $\nabla f = 0$ के पीछे छिपी रहती हैं। \emph{(ग) सदिश-मान वाली कोई
मध्यमान समता नहीं:} केवल \cref{thm:b2:diffcalc:mvi} की \emph{असमिका} बची रहती है (\cref{exo:b2:diffcalc:9}); $\R^m$,
$m \geq 2$ में मान लेने वाले $f$ के लिए $f(b) - f(a) = \dd
f_c(b-a)$ कभी मत लिखिए। \emph{(घ) स्थानीय
प्रतिलोमनीयता एकैकीपन नहीं है:} \cref{ex:b2:diffcalc:localnotglobal}। \emph{(ङ) प्रवणक अंतर्गुणन का है:}
$\nabla f$ किसी चुने हुए अंतर्गुणन में $\dd f_a$ को निरूपित करने वाला सदिश है; गुणन
बदलिए (जैसा द्विघात-रूप अध्याय के भारित उदाहरण में) और प्रवणक घूम जाता है,
जबकि \omterm{def:b2:diffcalc:differential}{अवकल} --- वह अंतर्निहित वस्तु --- टस से मस नहीं होता।
\end{remark}

\begin{remark}[यह कहाँ काम आता है]
इस अध्याय के आगे सब कुछ अवकल कलन का प्रयोग है: अवकल समीकरणों वाला अध्याय
प्रवाहों का रैखिकीकरण करता है और ल्यूविल का \omterm{def:b2:linalg:det}{सारणिक} सूत्र काम में लेता है
(जो इस अध्याय की सप्ताहांत समस्या में सिद्ध है); वक्रों और पृष्ठों वाले
अध्याय स्तर समुच्चयों तथा प्राचलनों का अध्ययन अंतर्निहित फलन प्रमेय के
द्वारा करते हैं; और बहु समाकल याकोबियों के द्वारा चर बदलते हैं। सप्ताहांत
समस्या स्वयं \emph{आव्यूहों की समष्टि पर} कलन विकसित करती है --- \omterm{def:b2:linalg:det}{सारणिक} का
\omterm{def:b2:diffcalc:differential}{अवकल}, प्रतिलोम का \omterm{def:b2:diffcalc:differential}{अवकल}, \omterm{ex:b2:nvs:matrixexp}{आव्यूह घातांकी}, और किसी चिकने स्तर समुच्चय के रूप
में लांबिक समूह --- अर्थात् जिसे तृतीय वर्ष का खंड बहुविध और ली समूहों के
रूप में औपचारिक बनाता है उसकी द्वितीय वर्ष वाली परछाईं।
\end{remark}

\section{अभ्यास}

\begin{exercise}[$\star$]\label{exo:b2:diffcalc:1}
$f(x,y) = (x^2 - y^2,\, 2xy)$ तथा $\Phi(r,\theta,z) = (r\cos\theta, r\sin\theta, z)$ के \omterm{def:b2:diffcalc:differential}{याकोबी आव्यूह} परिकलित कीजिए; \omterm{def:b2:diffcalc:differential}{अवकल} कहाँ प्रतिलोमनीय हैं?
\end{exercise}

\begin{exercise}[$\star$]\label{exo:b2:diffcalc:2}
$f(0,0) = 0$ के लिए $f(x,y) = \frac{x^3}{x^2 + y^2}$ लीजिए। सिद्ध कीजिए कि $0$ पर $f$ के सारे दिशिक अवकलज
विद्यमान हैं, परंतु $f$ $0$ पर \omterm{def:b2:diffcalc:differential}{अवकलनीय} नहीं है \emph{(प्रतिचित्रण $v \mapsto$
दिशिक अवकलज रैखिक नहीं है)}।
\end{exercise}

\begin{exercise}[$\star$]\label{exo:b2:diffcalc:3}
\cref{thm:b2:diffcalc:extrema} के द्वारा $f(x, y) = x^3 + y^3 -
3xy$ के, तथा $g(x,y) = x^4 +
y^4 - 2(x - y)^2$ के क्रांतिक बिंदु ज्ञात कीजिए और उनका
वर्गीकरण कीजिए।
\end{exercise}

\begin{exercise}[$\star\star$]\label{exo:b2:diffcalc:4}
मान लीजिए $f \colon \R^n \to \R$ $C^1$ है और \emph{घात $p$ का समघातीय}: $t > 0$ के लिए $f(tx) = t^pf(x)$। ऑयलर
की सर्वसमिका
\[
\langle \nabla f(x), x\rangle = p\,f(x) ,
\]
सिद्ध कीजिए, और $\R^n\setminus\{0\}$ पर $C^1$ फलनों के लिए उसका विलोम भी।
\end{exercise}

\begin{exercise}[$\star\star$]\label{exo:b2:diffcalc:5}
मान लीजिए $A$ \omterm{def:b2:quadratic:adjoint}{सममित} है और $f(x) = \frac12\langle Ax, x\rangle -
\langle b, x\rangle$। $\nabla f$ तथा $H_f$ परिकलित कीजिए; $f$ कब उत्तल
है? $A$ को धनात्मक निश्चित मानते हुए दिखाइए कि $f$ का $Ax = b$ के हल पर
अद्वितीय वैश्विक न्यूनतम है --- अर्थात् प्रवणक अवरोहण का औचित्य।
\end{exercise}

\begin{exercise}[$\star\star$]\label{exo:b2:diffcalc:6}
(लाग्रांज गुणक, एक प्रतिबंध, हाथ से सिद्ध) मान लीजिए $f, g$ $\R^2$ पर $C^1$ हैं,
और $f$ $\nabla g(a) \neq 0$ सहित स्तर समुच्चय $\{g = 0\}$ पर $a$ पर कोई स्थानीय चरम मान प्राप्त
करता है। सिद्ध कीजिए कि किसी $\lambda$ के लिए $\nabla f(a) = \lambda\nabla g(a)$। \emph{(अंतर्निहित फलन प्रमेय
से $a$ के पास स्तर समुच्चय का प्राचलन कीजिए और $t \mapsto f(\gamma(t))$ का अवकलन कीजिए।)}
अनुप्रयोग: वृत्त $x^2 + y^2 = 1$ पर $f(x,y) = xy$ के चरम मान।
\end{exercise}

\begin{exercise}[$\star\star$]\label{exo:b2:diffcalc:7}
$\R^2$ पर $f(x, y) = x^2 + y^2 - xy + x - y$ के चरम मान निर्धारित कीजिए, और फिर शीर्षों $(0,0)$, $(1,0)$, $(0,1)$ वाले
संवृत त्रिभुज पर उसका महत्तम और न्यूनतम \emph{(पहले अभ्यंतरीय क्रांतिक
बिंदु, फिर तीनों भुजाएँ, फिर शीर्ष)}।
\end{exercise}

\begin{exercise}[$\star\star\star$]\label{exo:b2:diffcalc:8}
$f \colon \R^2 \to \R^2$, $f(x, y) = (x + y^2,\; y + x^2)$ लीजिए। दिखाइए कि $f$ $0$ के पास स्थानीय अवकल समाकृतिकता है, $\dd(f^{-1})_{(0,0)}$
परिकलित कीजिए, और वह महत्तम $r$ ज्ञात कीजिए जिसके लिए $\dd f$ गोले $\norm{(x,y)}_2 < r$ पर
प्रतिलोमनीय हो \emph{($\det J_f$ परिकलित कीजिए)}।
\end{exercise}

\begin{exercise}[$\star\star\star$]\label{exo:b2:diffcalc:9}
(रोल विफल, मध्यमान बचा रहता है) ऐसा $f \colon \R \to \R^2$, $C^1$ दीजिए जिसके लिए $f(0) = f(2\pi)$ हो पर
सभी $t$ के लिए $f'(t) \neq 0$ (अर्थात् सदिश-मान वाला कोई रोल नहीं)। फिर अपने उदाहरण
पर \cref{thm:b2:diffcalc:mvi} की मध्यमान \emph{असमिका} सत्यापित कीजिए।
\end{exercise}

\begin{exercise}[$\star$]\label{exo:b2:diffcalc:10}
$\R^n$ पर $f(x) = \norm
x_2^2$ तथा $g(x) = \langle Ax, x\rangle$ के \omterm{def:b2:diffcalc:differential}{अवकल} और प्रवणक परिकलित कीजिए ($A$ कोई वर्ग आव्यूह,
\omterm{def:b2:quadratic:adjoint}{सममित} नहीं माना गया), और हर एक का हेस्सियन भी। किन $A$ के लिए $g$ उत्तल
है?
\end{exercise}

\begin{exercise}[$\star\star$]\label{exo:b2:diffcalc:11}
मान लीजिए $f \colon \R^n \to \R$ $C^2$ है। सिद्ध कीजिए कि $f$ उत्तल है यदि और केवल यदि उसका
हेस्सियन $H_x$ प्रत्येक $x$ पर धनात्मक अर्धनिश्चित हो \emph{(एक चर तक ले
आइए: $t \mapsto f(a + t(b-a))$; एक दिशा में टेलर--लाग्रांज लगाइए, और विलोम के लिए $\varphi''$ का
मूल्यांकन कीजिए)}।
\end{exercise}

\begin{exercise}[$\star\star\star$]\label{exo:b2:diffcalc:12}
(एक वैश्विक प्रतिलोम प्रमेय) मान लीजिए $g \colon \R^n \to \R^n$ $C^1$ है, सभी $x$ के लिए $\vertiii{\dd g_x} \leq k < 1$, और
$f =
\mathrm{id} + g$।
\begin{enumerate}
  \item दिखाइए $\norm{f(x) - f(y)} \geq (1 - k)\norm{x - y}$: अर्थात् $f$ एकैकी है और अपने प्रतिबिंब पर उसका
        प्रतिलोम \omterm{def:b2:metric:continuity}{संतत} है।
  \item दिखाइए कि प्रत्येक $y \in \R^n$ के लिए प्रतिचित्रण $x \mapsto y -
        g(x)$ \omterm{def:b2:metric:complete}{पूर्ण} समष्टि $\R^n$ का
        संकुचन है, और बानाख अचल बिंदु प्रमेय (\cref{thm:b2:metric:banach}) से निष्कर्ष निकालिए कि
        $f$ आच्छादक है।
  \item निष्कर्ष निकालिए कि $f$ $\R^n$ का एकैकी आच्छादक प्रतिचित्रण है जिसका
        प्रतिलोम $(1-k)^{-1}$-लिप्शिट्स है --- अर्थात् \cref{thm:b2:diffcalc:inverse} का \emph{वैश्विक}
        प्रतिरूप (जो इसके विपरीत विशुद्ध रूप से स्थानीय है)।
\end{enumerate}
\end{exercise}

\section{समस्या: आव्यूहों का कलन --- याकोबी, घातांकी और लांबिक समूह}

\begin{problem}\label{pb:b2:diffcalc:1}
अवकल कलन का सबसे स्वच्छ खेल-मैदान स्वयं समष्टि $\mathcal M_n(\R) \simeq \R^{n^2}$ है: उसके सबसे
स्वाभाविक प्रतिचित्रणों --- गुणनफल, प्रतिलोम, \omterm{def:b2:linalg:det}{सारणिक}, घातांकी --- के \omterm{def:b2:diffcalc:differential}{अवकल}
आश्चर्यजनक रूप से सुंदर हैं। यह समस्या उन सबका परिकलन करती है:
\emph{नॉयमान श्रेणी}, प्रतिलोम का \omterm{def:b2:diffcalc:differential}{अवकल}, \omterm{def:b2:linalg:det}{सारणिक} के लिए \emph{याकोबी
सूत्र}\index{याकोबी सूत्र} और लाभ में \emph{ल्यूविल का सूत्र}, $\det\eu^A = \eu^{\operatorname{tr}A}$ वाला
\emph{आव्यूह घातांकी}\index{आव्यूह घातांकी}, और अंत में लांबिक समूह $O_n$
किसी चिकने स्तर समुच्चय के रूप में जिसकी स्पर्श समष्टि प्रतिसममित आव्यूह
हैं --- अर्थात् भ्रूण-रूप में अवकल ज्यामिति। आगे सर्वत्र $\vertiii\cdot$ $\norm\cdot_2$ के अधीनस्थ
\omterm{thm:b2:nvs:continuouslinear}{संकारक मानदंड} है, और $\langle X, Y\rangle =
\operatorname{tr}(X^{\mathsf T}Y)$ फ्रोबेनियस अंतर्गुणन।

\medskip
\noindent\textbf{भाग I --- नॉयमान श्रेणी।}
\begin{enumerate}
  \item उपगुणात्मकता $\vertiii{AB} \leq
        \vertiii A\,\vertiii B$ सिद्ध कीजिए, और इससे निष्कर्ष निकालिए कि $A$ के
        बहुपद प्रतिचित्रण (आव्यूह गुणनफल, \omterm{def:b2:linalg:det}{सारणिक}, अनुरेख) $\mathcal M_n(\R)$ पर \omterm{def:b2:metric:continuity}{संतत}
        हैं।
  \item $\vertiii X < 1$ के लिए दिखाइए कि $\sum_{k\geq0}X^k$ $\mathcal M_n(\R)$ में \omterm{def:b2:series:def}{निरपेक्ष रूप से} अभिसरित होती है
        (\cref{thm:b2:nvs:absoluteconvergence}), कि उसका योग $(I - X)^{-1}$ है, और यह कि
        \[
        \vertiii{(I - X)^{-1}} \leq \frac{1}{1 - \vertiii X},
        \qquad
        (I - X)^{-1} = I + X + O\bigl(\vertiii X^2\bigr) .
        \]
  \item इससे निष्कर्ष निकालिए कि $GL_n(\R)$ \emph{\omterm{def:b2:metric:topology}{विवृत}} है: यदि $A$ प्रतिलोमनीय
        हो और $\vertiii H <
        \frac{1}{\vertiii{A^{-1}}}$, तो $A + H$ प्रतिलोमनीय है। यह भी निष्कर्ष निकालिए कि $GL_n(\R)$
        $\mathcal M_n(\R)$ में \emph{सघन} है \emph{($A$ को $\varepsilon I$ से विक्षुब्ध कीजिए:
        $\det(A + \varepsilon I)$ $\varepsilon$ में अशून्य बहुपद है)}।
  \item दिखाइए कि प्रतिलोमन प्रतिचित्रण $\Phi(A) = A^{-1}$ $GL_n(\R)$ पर \omterm{def:b2:metric:continuity}{संतत} है।
\end{enumerate}

\medskip
\noindent\textbf{भाग II --- पहले \omterm{def:b2:diffcalc:differential}{अवकल}।}
\begin{enumerate}[resume]
  \item दिखाइए कि वर्ग-प्रतिचित्रण $A \mapsto A^2$ \omterm{def:b2:diffcalc:differential}{अवकलनीय} है और उसका \omterm{def:b2:diffcalc:differential}{अवकल} $H \mapsto AH + HA$ है,
        और अधिक सामान्य रूप में $A \mapsto A^k$ का \omterm{def:b2:diffcalc:differential}{अवकल} $H \mapsto \sum_{i=0}^{k-1}
        A^iHA^{k-1-i}$। सामान्यतः $kA^{k-1}H$ क्यों
        \emph{नहीं} लिखा जा सकता?
  \item सिद्ध कीजिए कि $\Phi(A) = A^{-1}$ $GL_n(\R)$ पर \omterm{def:b2:diffcalc:differential}{अवकलनीय} है और
        \[
        \dd\Phi_A(H) = -A^{-1}HA^{-1}
        \]
        \emph{($(A + H)^{-1} = (I + A^{-1}H)^{-1}A^{-1}$ लिखिए और प्रश्न 2 से प्रसार कीजिए)}। सूत्र को अदिश
        स्थिति $n = 1$ के विरुद्ध जाँचिए।
  \item किसी $C^1$ वक्र $t \mapsto A(t) \in GL_n(\R)$ के लिए $\bigl(A(t)^{-1}\bigr)' = -A^{-1}A'A^{-1}$ निष्कर्ष रूप में निकालिए, और $t = 0$
        पर $t \mapsto (I + tB)^{-1}$ का प्रथम कोटि तक प्रसार कीजिए।
  \item $f(A) =
        \operatorname{tr}(A^k)$ का \omterm{def:b2:diffcalc:differential}{अवकल} परिकलित कीजिए और फ्रोबेनियस अंतर्गुणन के लिए उसका
        प्रवणक पहचानिए:
        \[
        \dd f_A(H) = k\operatorname{tr}\bigl(A^{k-1}H\bigr),
        \qquad
        \nabla f(A) = k\,\bigl(A^{k-1}\bigr)^{\mathsf T} .
        \]
  \item $f(A) = \operatorname{tr}
        (A^{\mathsf T}A) = \norm A_F^2$ के लिए वही प्रश्न: \omterm{def:b2:diffcalc:differential}{अवकल}, प्रवणक, और (अचर) हेस्सियन; निष्कर्ष
        निकालिए कि $\norm\cdot_F^2$ यथार्थतः उत्तल है।
\end{enumerate}

\medskip
\noindent\textbf{भाग III --- याकोबी सूत्र।}
\begin{enumerate}[resume]
  \item सिद्ध कीजिए
        \[
        \det(I + H) = 1 + \operatorname{tr}H +
        O\bigl(\vertiii H^2\bigr)
        \]
        \emph{($\det(e_1 + h_1, \dots, e_n + h_n)$ को स्तंभों में बहुरैखिकता से खोलिए: कम से कम दो
        $h$-स्तंभों वाले पद $O(\vertiii H^2)$ हैं)}: $\dd(\det)_I =
        \operatorname{tr}$।
  \item प्रतिलोमनीय $A$ के लिए निष्कर्ष निकालिए
        \[
        \dd(\det)_A(H) = \det(A)\,
        \operatorname{tr}\bigl(A^{-1}H\bigr) .
        \]
  \item दिखाइए कि \emph{प्रत्येक} $A$ के लिए (प्रतिलोमनीय हो या न हो)
        $\frac{\partial\det}{\partial a_{ij}}(A) = C_{ij}$ $(i,j)$ सहगुणनखंड है \emph{(पंक्ति $i$ के अनुदिश लाप्लास प्रसार)},
        अतः संलग्नक $\operatorname{adj}A =
        \operatorname{com}(A)^{\mathsf T}$ के साथ:
        \[
        \dd(\det)_A(H) =
        \operatorname{tr}\bigl(\operatorname{adj}(A)\,H\bigr),
        \qquad
        \nabla(\det)(A) = \operatorname{com}(A) ,
        \]
        जो $A$ के प्रतिलोमनीय होने पर प्रश्न 11 पुनः दे देता है ($\operatorname{adj}A = \det(A)A^{-1}$)।
        यही \emph{याकोबी सूत्र} है: $\bigl(\det
        A(t)\bigr)' = \operatorname{tr}\bigl(
        \operatorname{adj}(A(t))\,A'(t)\bigr)$।
  \item (ल्यूविल का सूत्र) मान लीजिए $A(t)$ आव्यूहों का ऐसा $C^1$ वक्र है जो
        रैखिक अवकल समीकरण $A'(t) = M(t)A(t)$ पूरा करता है। सिद्ध कीजिए
        \[
        \bigl(\det A(t)\bigr)' =
        \operatorname{tr}\bigl(M(t)\bigr)\,\det A(t),
        \qquad\text{अतः}\qquad
        \det A(t) = \det A(0)\,
        \exp\Bigl(\int_0^t\operatorname{tr}M\Bigr)
        \]
        \emph{($\operatorname{adj}(A)\,A = \det(A)I$ तथा अनुरेख की \omterm{def:b2:structures:generated}{चक्रीय} अपरिवर्त्यता का उपयोग कीजिए)}
        --- यही वह व्रोंस्कियन सर्वसमिका है जिसे अवकल समीकरणों वाला अध्याय
        लगातार काम में लेगा।
  \item दिखाइए कि $SL_n(\R) = \{\det = 1\}$ चिकना स्तर समुच्चय है: प्रत्येक $A \in SL_n(\R)$ पर \omterm{def:b2:diffcalc:differential}{अवकल} $\dd(\det)_A$
        $\R$ पर \emph{आच्छादक} रैखिक प्रतिचित्रण है \emph{(उसका $H = \frac1nA$ पर
        मूल्यांकन कीजिए)}।
\end{enumerate}

\medskip
\noindent\textbf{भाग IV --- \omterm{ex:b2:nvs:matrixexp}{आव्यूह घातांकी}।}
\begin{enumerate}[resume]
  \item दिखाइए कि $\eu^A = \sum_{k\geq0}\frac{A^k}{k!}$ प्रत्येक $A$ के लिए \omterm{def:b2:series:def}{निरपेक्ष रूप से} और हर गोले पर
        \omterm{def:b2:funcseq:series}{सामान्य रूप से} अभिसरित होती है, जहाँ $\vertiii{\eu^A} \leq
        \eu^{\vertiii A}$; और यह कि $\eu^A$ $A$ पर
        \omterm{def:b2:metric:continuity}{संतत} रूप से निर्भर करता है।
  \item सिद्ध कीजिए कि $AB = BA$ से $\eu^{A+B} =
        \eu^A\eu^B$ निकलता है \emph{(\omterm{thm:b2:series:fubini}{कोशी गुणनफल}, जो
        निरपेक्ष अभिसरण से वैध है)}; और इससे निष्कर्ष निकालिए कि $\eu^A$ सदा
        प्रतिलोमनीय है जिसका प्रतिलोम $\eu^{-A}$ है: अर्थात् $\exp$ $\mathcal M_n(\R)$ को $GL_n(\R)$ में
        भेजता है।
  \item दिखाइए कि $t \mapsto \eu^{tA}$ $C^1$ है (वस्तुतः $C^\infty$) और
        \[
        \frac{\dd}{\dd t}\,\eu^{tA} = A\,\eu^{tA} =
        \eu^{tA}A
        \]
        \emph{(खंडों पर श्रेणी का पद-दर-पद अवकलन कीजिए: व्युत्पन्न श्रेणी
        \omterm{def:b2:funcseq:series}{सामान्य रूप से} अभिसरित होती है)}।
  \item सर्वसमिका
        \[
        \det\bigl(\eu^{A}\bigr) = \eu^{\operatorname{tr}A}
        \]
        सिद्ध कीजिए \emph{($A(t) = \eu^{tA}$ पर ल्यूविल का सूत्र, प्रश्न 13, लगाइए)}।
        तर्कसंगति की जाँचें: $n = 1$; शून्यंभावी $A$; और अनुरेख-शून्य
        आव्यूह $SL_n(\R)$ में उतरते हैं।
  \item दिखाइए कि $\eu^H = I + H + O(\vertiii H^2)$, अतः $\exp$ $0$ पर $\dd(\exp)_0 =
        \mathrm{id}$ के साथ \omterm{def:b2:diffcalc:differential}{अवकलनीय} है; और
        प्रतिलोम फलन प्रमेय (\cref{thm:b2:diffcalc:inverse}) से निष्कर्ष निकालिए कि $\exp$ $0$ के
        किसी प्रतिवेश से $I$ के किसी प्रतिवेश पर $C^1$ अवकल समाकृतिकता
        है: अर्थात् तत्समक के निकट हर आव्यूह का लघुगणक होता है।
  \item दिखाइए कि $\exp$ \omterm{def:b2:quadratic:adjoint}{सममित} आव्यूहों को \omterm{def:b2:quadratic:adjoint}{सममित} \emph{धनात्मक निश्चित}
        आव्यूहों पर एकैकी आच्छादक रूप से भेजता है \emph{(विकर्णन कीजिए;
        प्रतिलोम वर्णक्रमीय लघुगणक है)}।
\end{enumerate}

\medskip
\noindent\textbf{भाग V --- स्तर समुच्चय के रूप में लांबिक समूह।}
\begin{enumerate}[resume]
  \item $\mathcal M_n(\R)$ से \omterm{def:b2:quadratic:adjoint}{सममित} आव्यूहों $S_n$ में जाने वाला $F(A) = A^{\mathsf T}A$ लीजिए। $\dd F_A(H) =
        A^{\mathsf T}H + H^{\mathsf T}A$ परिकलित
        कीजिए और दिखाइए कि प्रत्येक $A \in O_n = F^{-1}(I)$ पर यह \omterm{def:b2:diffcalc:differential}{अवकल} $S_n$ पर \emph{आच्छादक}
        है \emph{($S \in S_n$ दिया हो तो $H = \frac12 AS$ आज़माइए)}: अर्थात् $O_n$ विमा $n^2 - \frac{n(n+1)}2 = \frac{n(n-1)}2$ का
        चिकना स्तर समुच्चय है।
  \item दिखाइए कि $A(0) =
        I$ वाले हर $C^1$ वक्र $A(t) \in O_n$ का वेग $A'(0)$
        \emph{प्रतिसममित} होता है, और विलोमतः प्रतिसममित $K$ के लिए वक्र
        $\eu^{tK}$ $O_n$ में बना रहता है: अर्थात् $I$ पर $O_n$ की स्पर्श समष्टि
        ठीक प्रतिसममित आव्यूह हैं।
  \item प्रतिसममित $K$ के लिए $\det\eu^{K} = 1$ दिखाइए (प्रश्न 18): अतः घातांकी वक्र
        घूर्णन समूह $SO_n$ में बसता है। $n = 2$ के लिए उसे पूरा परिकलित कीजिए:
        $J =
        \begin{pmatrix}0 & -1\\ 1 & 0\end{pmatrix}$ के साथ सिद्ध कीजिए
        \[
        \eu^{\theta J} = \begin{pmatrix} \cos\theta &
        -\sin\theta\\ \sin\theta & \cos\theta\end{pmatrix} :
        \]
        अर्थात् \omterm{ex:b2:nvs:matrixexp}{आव्यूह घातांकी} $\theta$ से घूर्णन \emph{ही} है, और कोसाइन तथा
        साइन की श्रेणी-परिभाषाएँ किसी आव्यूह के भीतर फिर प्रकट हो जाती
        हैं।
  \item (सघनता की युक्ति) $GL_n(\R)$ की सघनता (प्रश्न 3) और संततता का उपयोग
        करके सर्वसमिका
        \[
        \operatorname{adj}(AB) =
        \operatorname{adj}(B)\operatorname{adj}(A)
        \]
        को प्रतिलोमनीय से सभी आव्यूहों तक बढ़ाइए \emph{(प्रतिलोमनीय $A, B$
        के लिए दोनों पक्ष $\det(AB)(AB)^{-1}$ हैं; और दोनों पक्ष प्रविष्टियों में बहुपद
        हैं)}।
  \item संश्लेषण। एक-एक वाक्य में: (क) हर भाग का इंजन किन पिछले अध्यायों
        ने दिया (पूर्णता और मानदंडित बीजगणित; वर्णक्रमीय प्रमेय; प्रतिलोम
        फलन प्रमेय); (ख) इस समस्या का कौन-सा सूत्र अवकल समीकरणों वाला
        अध्याय काम में लेगा, और कहाँ; (ग) $\dd(\det)_I = \operatorname{tr}$ और $\det\eu^A =
        \eu^{\operatorname{tr}A}$ अनुरेख तथा \omterm{def:b2:linalg:det}{सारणिक}
        के बारे में ``अत्यल्प और वैश्विक आयतन'' के रूप में क्या कहते हैं;
        (घ) तृतीय वर्ष का खंड प्रश्न 21--23 का क्या बना देता है (ली समूह
        और उनके ली बीजगणित)।
\end{enumerate}
\end{problem}
